\documentclass[10pt,a4paper]{article}
\usepackage[a4paper,margin=0.72in]{geometry}
\usepackage[utf8]{inputenc}
\usepackage[T1]{fontenc}
\usepackage{lmodern}
\usepackage{amsmath,amssymb}
\usepackage{booktabs}
\usepackage{array}
\usepackage{graphicx}
\usepackage{float}
\usepackage{enumitem}
\usepackage{hyperref}
\usepackage{xurl}
\usepackage{fontspec}
\newfontfamily\devanagarifont{Kohinoor Devanagari}[Script=Devanagari]
\newcommand{\deva}[1]{{\devanagarifont #1}}

\hypersetup{colorlinks=true,linkcolor=blue,urlcolor=blue,citecolor=blue}
\setlist[itemize]{noitemsep,topsep=2pt}
\setlist[enumerate]{noitemsep,topsep=2pt}

\title{\textbf{Language Models and Agents}\\\large Individual Project Report}
\author{2023102073\\Harshita Kumari}
\date{September 2026}

\begin{document}
\maketitle

\begin{abstract}
Two independent approximately 25M-parameter decoder-only Transformer language models were built and trained entirely from scratch for Hindi and Nepali. The models use separate corpora, tokenizers, and weights. Approximately 500M tokens were used for pretraining in each language. The models were evaluated using perplexity, bits-per-byte, generation metrics, and attention analysis, followed by finetuning on a synthetic reasoning task. Hindi obtained a test perplexity of 31.96 and 36.75\% reasoning accuracy after finetuning, while Nepali obtained a test perplexity of 40.33 and 8.75\% reasoning accuracy.
\end{abstract}

\section{Phase 1: Data Collection and Tokenizer}

\subsection{Data Collection}
The corpus was built using a combination of FineWeb2, IndicCorp, Sangraha,
Wikipedia, and several manually collected web sources. For the manual
collection, we used sitemaps and archives to find relevant pages and
applied rate limits at the domain level. Sources that explicitly
restricted AI-related crawling were not included.

\begin{table}[H]\centering
\caption{Final corpus statistics}
\small
\begin{tabular}{lrr}\toprule
Metric & Hindi & Nepali\\\midrule
Total tokens & 499,999,997 & 499,302,601\\
Manual token fraction & 20.29\% & 19.89\%\\
Manual tokens & 101,450,874 & 99,302,601\\
Downloaded tokens & 398,549,123 & 400,000,000\\
\bottomrule\end{tabular}
\end{table}

The Nepali corpus fell slightly short of the 500M-token target, by
697,399 tokens (0.14\%).



\subsection{Data Processing Pipeline}
Before training the tokenizers, both corpora were processed using the
same four-stage pipeline, with the entire process carried out
independently for each language.

\paragraph{1. Cleaning, deduplication, and sampling.}
The data was processed in this order:

\begin{enumerate}
\item \textbf{Unicode normalization and mojibake correction} using
\texttt{ftfy}. This helped fix encoding issues that can occur in
scraped or downloaded text. NFC normalization ensured that
visually identical characters were represented in the same way before
duplicate removal.

\item \textbf{Script filtering.} We kept lines containing at least
60\% Devanagari characters to remove irrelevant non-Devanagari
boilerplate while still keeping naturally occurring mixed-script
sentences.

\item \textbf{Whitespace and punctuation normalization.} Removed extra
whitespace and inconsistent punctuation to make text more uniform.

\item \textbf{Exact duplicate removal.} We used SHA-256 hashes of the
normalized, lowercased, and whitespace-collapsed text to identify
duplicates. Manual documents were given priority when a manual and
downloaded document had the same hash.

\item \textbf{Minimum-length filtering.} Documents containing fewer
than 20 characters after cleaning were removed.

\item \textbf{Sampling to the target token count.} After cleaning, the
remaining documents were sampled until the required token target was
reached, while following the specified manual-to-downloaded data ratio.
\end{enumerate}

\paragraph{2. Splitting.}
The data was split at the \textbf{document level} rather than splitting
individual sentences or portions of documents. This prevents parts of
the same document from appearing in both the training and evaluation
sets and helps avoid data leakage.

\begin{table}[H]\centering
\caption{Train/validation/test split sizes (document counts)}
\small
\begin{tabular}{lrrr}\toprule
Language & Train & Val & Test \\\midrule
Hindi & 1,340,887 (98.00\%) & 13,682 (1.00\%) & 13,684 (1.00\%)\\
Nepali & 1,855,143 (98.00\%) & 18,930 (1.00\%) & 18,931 (1.00\%)\\
\bottomrule\end{tabular}
\end{table}

The pipeline is summarized as:
\[
{raw} \;\rightarrow\; {clean/dedup/filter/sample}
\;\rightarrow\; {split (train/val/test)}
\;\rightarrow\; {train  tokenizer (train  split  only)}.
\]

\subsection{Tokenizer Training}
\paragraph{Algorithm: SentencePiece BPE.}
We used Byte-Pair Encoding (BPE) to build the tokenizer. It starts with
individual characters as the basic vocabulary and gradually combines the
most frequently occurring adjacent pieces into larger pieces. This process
continues until the vocabulary reaches the target size of 8,000 tokens.
As a result, frequently occurring words or parts of words can be represented
using a single token, while less common words are broken down into smaller
pieces. This also helps keep the unknown-token rate low, since unseen words
can still be represented using smaller subword pieces or individual
characters.

We used SentencePiece because it does not require the text to be split into
words before tokenization. This is useful for Devanagari text, where
whitespace does not always correspond perfectly to word boundaries.

\begin{table}[H]\centering
\caption{Final tokenizer statistics}
\small
\begin{tabular}{lrr}\toprule
Metric & Hindi & Nepali\\\midrule
Vocabulary size & 8,000 & 8,000\\
Fertility (tokens/word) & 1.380 & 1.536\\
Unknown-token rate & 0.261\% & 0.190\%\\
\bottomrule\end{tabular}
\end{table}

\paragraph{Choosing the vocabulary size.}
To decide the vocabulary size, we tested different values ranging from
2,000 to 32,000 on a validation sample. The unknownb token rate remained below
0.33\% for both languages across all vocabulary sizes, so this was not a
major factor in our decision.

The main trade off was between fertility and the number of parameters
required by the embedding table. Since the embedding size grows linearly
with the vocabulary size ($\text{vocab\_size}\times d_{\text{model}}$),
using a much larger vocabulary increases the parameter cost significantly.
For example, increasing the vocabulary from 8K to 32K improves Hindi's
fertility by only about 13\% (1.38 to 1.20 tokens/word), while the
embedding table increases from 16.38\% to 65.54\% of the 25M-parameter
budget. Based on this trade-off, we selected a vocabulary size of 8,000,
where the improvement from using a larger vocabulary was no longer worth
the additional parameter cost.

\begin{table}[H]\centering
\caption{Vocab-size sweep}
\label{tab:vocab-sweep}
\small
\begin{tabular}{rrrrrr}\toprule
& \multicolumn{2}{c}{Hindi} & \multicolumn{2}{c}{Nepali} & \\
Vocab size & Fertility & UNK\% & Fertility & UNK\% & Emb.\ \% of 25M\\\midrule
2,000  & 1.7747 & 0.218\% & 2.1043 & 0.139\% & 4.10\%\\
4,000  & 1.5455 & 0.250\% & 1.7778 & 0.164\% & 8.19\%\\
\textbf{8,000}  & \textbf{1.3795} & \textbf{0.280\%} & \textbf{1.5364} & \textbf{0.190\%} & \textbf{16.38\%}\\
16,000 & 1.2667 & 0.305\% & 1.3728 & 0.213\% & 32.77\%\\
24,000 & 1.2219 & 0.316\% & 1.3022 & 0.224\% & 49.15\%\\
32,000 & 1.1988 & 0.322\% & 1.2621 & 0.231\% & 65.54\%\\
\bottomrule\end{tabular}
\end{table}

\begin{figure}[H]\centering
\includegraphics[width=\textwidth]{vocab_size_sweep.png}
\caption{Fertility (tokens/word), UNK rate, and embedding cost vs.\ vocab size (8,000 marked).}
\label{fig:vocab-sweep}
\end{figure}

\section{Phase 2: Model, Pretraining and Evaluation}

\subsection{Architecture}
Both models use the same architecture while
keeping their weights fully independent.

\begin{table}[H]\centering
\caption{Model configuration}
\label{tab:model-config}
\small
\begin{tabular}{lr}\toprule
Parameter & Value\\\midrule
Vocabulary & 8,000\\
Model dimension $d_{\text{model}}$ & 512\\
Transformer layers & 7\\
Attention heads & 8\\
Head dimension $d_k$ & 64\\
FFN dimension & 2,048\\
Context length & 512\\
Dropout & 0.1\\
Output head & Tied to embeddings\\
Total parameters & 26,425,856\\
\bottomrule\end{tabular}
\end{table}

\paragraph{Positional embeddings.} Learned absolute positional embeddings
are used instead of sinusoidal or RoPE: at this scale, with a fixed 512-token
context, a learned $512\times512$ table costs only $\sim$1\% of the model
and is simpler to verify, at the cost of no length extrapolation beyond 512
tokens.

\paragraph{Causal self-attention.} Implemented from four independent
$W_Q,W_K,W_V,W_O$ projections (not a fused QKV matrix, for readability):
\[
\operatorname{Attention}(Q,K,V)=
\operatorname{softmax}\!\left(\frac{QK^T}{\sqrt{d_k}}+M\right)V,
\]
with $h=8$ heads of $d_k=512/8=64$ each. The
$1/\sqrt{d_k}$ scale keeps pre-softmax logit variance constant regardless of
head size -- without it, variance grows linearly with $d_k$ and pushes
softmax toward a near-one-hot regime with vanishing gradients. $M$ is an
additive causal mask ($0$ / $-\infty$), and its correctness was verified
empirically (perturbing only the last token of a random sequence and
confirming every earlier logit is unchanged): max logit difference was
exactly $0.0$ across 5 trials.

\paragraph{Transformer block.} Pre-norm self-attention (residual) $\to$
pre-norm feed-forward (residual). Pre-norm was chosen over post-norm for
training stability -- post-norm GPT-style models are markedly more sensitive
to learning-rate/warmup, a risk not worth taking on a limited,
hard-to-diagnose Colab compute budget. The FFN uses GELU with inner
dimension $4\times d_{\text{model}}=2048$.

\paragraph{Weight tying.} The output head is tied to the token embedding
matrix, saving an entire $8000\times512=4{,}096{,}000$-parameter matrix --
13.4\% of the model (30.5M untied vs.\ 26.4M tied), a meaningful fraction at
this parameter budget.

\begin{table}[H]\centering
\caption{Exact parameter breakdown (sums to the reported total)}
\label{tab:param-breakdown}
\small
\begin{tabular}{lr}\toprule
Component & Parameters\\\midrule
Token embedding (tied with output head) & 4,096,000\\
Positional embedding & 262,144\\
7 $\times$ Transformer block (3,152,384 each) & 22,066,688\\
Final LayerNorm & 1,024\\
\midrule
\textbf{Total} & \textbf{26,425,856}\\
\bottomrule\end{tabular}
\end{table}

\paragraph{Depth vs.\ width.} With an 8,000-token vocabulary, the embedding
tables cost a near-fixed $\sim$4.4M, leaving $\sim$22M for the Transformer
stack. A narrower-but-deeper configuration (512 dim, 7 layers) was chosen
over wider-shallower because depth composes hierarchical/longer-range
structure more effectively per parameter at this scale, while keeping
per-step compute tractable for a Colab session. $d_{\text{model}}=512$ and
$n_{\text{heads}}=8$ were kept as well-precedented defaults rather than
swept, since architecture search was out of scope.




\subsection{Pretraining Procedure}
Rather than going through the corpus sequentially from beginning to end,
we used random sampling during training. At each step, 32 random
512-token windows were selected from the tokenized corpus, with
replacement. This means that the same part of the corpus could be sampled
more than once, while some parts might not be seen during the training run.

Gradients were accumulated over 4 micro-batches before updating the model.
Therefore, each optimizer step processed

\[
32 \times 4 \times 512 = 65{,}536 \text{ tokens}.
\]

Training was carried out for a fixed total of 8,000 steps.

\paragraph{Why random sampling with replacement instead of epochs.}
We used a fixed token budget of approximately 500M tokens rather than
defining training in terms of epochs. In fact, the Hindi corpus alone
contains around 675M tokenized tokens, which is already larger than the
524M-token training budget. Therefore, training for a complete epoch would
not be possible within the given budget.

This left us with two possible approaches: making one deterministic pass
through a shuffled version of the corpus and stopping when the token budget
was reached, or randomly sampling windows with replacement. We chose the
second approach because it does not require creating and storing a
precomputed shuffle or index for the entire corpus. It also makes it easy
to resume training after an interruption using just the training step
count. This was particularly useful for the Nepali model, whose training
had to be restarted four times. Random sampling with replacement is also
consistent with the approach commonly used in large-corpus language model
training, such as GPT-2, GPT-3, and nanoGPT, when the training budget is
smaller than the available corpus.



\begin{table}[H]\centering
\caption{Pretraining hyperparameters}
\small
\begin{tabular}{ll}\toprule
Setting & Value\\\midrule
Optimizer & AdamW, $\beta=(0.9,0.95)$\\
Peak learning rate & $3\times10^{-4}$\\
Schedule & linear warmup (200 steps) $\to$ cosine decay to 10\% of peak\\
Weight decay & 0.1, decoupled, $\geq$2D matrices only\\
Gradient clipping & max norm 1.0\\
Batch $\times$ accumulation $\times$ block & $32\times4\times512=65{,}536$ tokens/step\\
Total steps & 8,000 ($\approx$524M tokens)\\
Checkpointing & every 200 steps, resumable (\texttt{latest.pt} + \texttt{best.pt})\\
\bottomrule\end{tabular}
\end{table}

\begin{table}[H]\centering
\caption{Pretraining results}
\small
\begin{tabular}{lrr}\toprule
Metric & Hindi & Nepali\\\midrule
Tokens seen & 524,288,000 & 524,288,000\\
Final train loss & 3.5647 & 3.7242\\
Final validation loss & 3.4717 & 3.6646\\
Validation perplexity & \textbf{32.19} & \textbf{39.04}\\
\bottomrule\end{tabular}
\end{table}


\subsection{Intrinsic Evaluation}
Perplexity was calculated as
\[
PPL=e^{\overline{\ell}},
\]
where $\overline{\ell}$ is mean cross-entropy in nats per token.
Bits-per-byte was calculated as
\[
BPB=\frac{\sum_i NLL_i}
{\ln(2)\times\text{UTF-8 bytes}}.
\]

\begin{table}[H]\centering
\caption{Test-set intrinsic evaluation}
\small
\begin{tabular}{lrr}\toprule
Metric & Hindi & Nepali\\\midrule
Test tokens & 6,630,400 & 7,754,752\\
Cross-entropy & 3.4646 & 3.6971\\
Perplexity & \textbf{31.96} & \textbf{40.33}\\
Bits-per-byte & \textbf{0.5259} & \textbf{0.4810}\\
\bottomrule\end{tabular}
\end{table}

Since the Hindi and Nepali tokenizers split text into different numbers of tokens, their token-level perplexities are not directly comparable. BPB provides an additional measure that evaluates performance at the byte level.

\subsection{Generation Evaluation}
Generation was evaluated using held-out prefixes with greedy decoding
and temperatures 0.5, 1.0, and 1.5. BLEU-4, chrF++, ROUGE-L,
Distinct-1, and repetition were measured.

\begin{table}[H]\centering
\caption{Generation metrics}
\scriptsize
\begin{tabular}{llrrrrr}\toprule
Model & Setting & BLEU & chrF++ & ROUGE-L & D-1 & Rep.\\\midrule
H & Greedy & 1.76 & 13.88 & .255 & .132 & .603\\
H & $T=0.5$ & 1.49 & 14.91 & .247 & .223 & .222\\
H & $T=1.0$ & .86 & 16.52 & .252 & .463 & .003\\
H & $T=1.5$ & .12 & 14.26 & .233 & .676 & .000\\
L & Greedy & 1.75 & 14.60 & .212 & .176 & .557\\
L & $T=0.5$ & 1.91 & 16.90 & .242 & .336 & .169\\
L & $T=1.0$ & .84 & 18.30 & .233 & .598 & .002\\
L & $T=1.5$ & .10 & 15.97 & .228 & .754 & .000\\
\bottomrule\end{tabular}
\end{table}

Greedy decoding showed substantial repetition. Higher temperature
reduced repetition and increased diversity, but very high temperature
also reduced reference overlap.

\subsection{Attention Analysis}
Attention entropy and mean attention distance were measured for every
(layer, head) pair, on 200 held-out test sentences (up to 128 tokens each):
\[
H_{l,h}=-\frac{1}{T}\sum_t\sum_k
p_{t,k}^{(l,h)}\log p_{t,k}^{(l,h)},
\qquad
D_{l,h}=\frac{1}{T}\sum_t\sum_k
p_{t,k}^{(l,h)}\,|t-k|,
\]
where $p^{(l,h)}$ is the softmax attention weight matrix for layer $l$, head
$h$. Lower entropy means more peaked (focused) attention; larger distance
means the head attends further from the query position.

\begin{table}[H]\centering
\caption{Mean attention entropy and distance by layer (averaged over 8 heads)}
\small
\begin{tabular}{lrrrr}\toprule
& \multicolumn{2}{c}{Entropy (nats)} & \multicolumn{2}{c}{Mean distance (tokens)}\\
Layer & Hindi & Nepali & Hindi & Nepali\\\midrule
0 (early) & 3.29 & 3.31 & 20.39 & 20.51\\
1 & 3.36 & 3.36 & 23.40 & 24.13\\
2 & 2.92 & 3.11 & 14.81 & 17.94\\
3 & 2.47 & 2.30 & 8.67 & 6.46\\
4 (most local) & 1.74 & 1.60 & 6.22 & 5.11\\
5 & 2.08 & 2.19 & 13.00 & 15.34\\
6 (late) & 2.20 & 2.20 & \textbf{31.29} & \textbf{31.58}\\
\bottomrule\end{tabular}
\end{table}

\begin{figure}[H]\centering
\includegraphics[width=1\textwidth]{hindi_attn_ex0_earlylayer0.png}\\
\includegraphics[width=1\textwidth]{hindi_attn_ex0_latelayer6.png}
\caption{Model H (Hindi): layer 0 (early, diffuse) vs.\ layer 6 (late, sink-dominated) attention heatmaps, heads 0--3.}
\end{figure}
\paragraph{Early layers (0--1): broad and less specialized attention.}
In the first few layers, the attention entropy remains relatively high,
around 3.3--3.4 nats, while the mean attention distance is also fairly
large ($\sim$20--24 tokens). This suggests that the heads are not yet
strongly specialized and tend to spread their attention across a wider
part of the context. At this stage, the model has not built up enough
higher-level representations to focus strongly on particular positions.

\paragraph{Middle layers (3--4): more local and positional attention.}
In the middle layers, the attention patterns become much more focused.
Entropy drops to as low as 0.19--0.75 nats for some heads, and the mean
attention distance decreases to around 5--9 tokens. This indicates that
some heads start focusing mainly on nearby tokens and positional
relationships. For Hindi and Nepali, these local dependencies can include
adjacent-word relationships and attached postpositions or case markers.
Such short-range patterns are useful for predicting the next token, so
the model appears to dedicate some of its attention capacity to these
local relationships.

\paragraph{Layer 6 (late): large distances caused mainly by an attention
sink.}
In the final layer, the mean attention distance increases again to around
31 tokens, which is the highest among the layers. However, this does not
necessarily mean that the model is performing more long-range reasoning.
The heatmaps suggest that the increase is mainly caused by an
\textbf{attention sink}. Several heads put a large amount of their
attention on a fixed early token, usually position 0, from many different
query positions.

This can happen because the attention weights have to sum to 1 across
all positions. When a head does not have a particularly relevant token
to focus on, it may place much of its attention on a fixed token that is
always present. As a result, the average distance becomes larger even
though the head is not necessarily using meaningful long-range
information. This is why the heatmaps are important for interpreting
the results rather than relying only on the entropy and distance values.
At the same time, a few heads in this layer do show genuine
long-range, content-based attention rather than the attention-sink
pattern.



\section{Phase 3: Reasoning Finetuning}

\subsection{Synthetic Dataset}
Each language's reasoning finetune set was generated programmatically by a
shared, language-agnostic generator (\texttt{reasoning\_gen.py}), with
per-language word lists and templates supplying all actual text -- so the
two languages share generation \emph{logic} but no content. Three task
types are covered:
\begin{itemize}
\item \textbf{Direct}: compare 2 entities on one attribute (age, height,
price, or quantity) -- answer is the greater, smaller, or ``both equal''.
\item \textbf{Transitive}: 3 entities chained $A\sim B\sim C$ -- answer is
the overall extreme (min/max).
\item \textbf{Multihop}: 3 entities chained, but the question asks about
the \emph{derived}, non-adjacent relation between $A$ and $C$ -- the only
task type that requires actually chaining both stated facts rather than
reading off an implied order.
\end{itemize}

\begin{table}[H]\centering
\caption{Reasoning dataset statistics (identical for Hindi and Nepali)}
\small
\begin{tabular}{lrrr}\toprule
& Train & Val & Test\\\midrule
Total examples & 6,000 & 800 & 800\\
\quad Direct & 2,024 & 277 & 368\\
\quad Transitive & 1,991 & 258 & 227\\
\quad Multihop & 1,985 & 265 & 205\\\midrule
\quad Age & 1,515 & 216 & 260\\
\quad Height & 1,519 & 189 & 263\\
\quad Price & 1,471 & 199 & 156\\
\quad Quantity & 1,495 & 196 & 121\\
\bottomrule\end{tabular}
\end{table}

\paragraph{Leakage control.} Both entity names \emph{and} sentence
templates were split into disjoint train/test pools before generation, so
every test example uses an entity name \emph{and} a template phrasing the
model never saw during finetuning -- 6 held-out person names and 3 held-out
object names per language, with at least one template per task type
reserved for test only.

% --- Option A: native script (requires XeLaTeX/LuaLaTeX + fontspec/polyglossia + a Devanagari font) ---
\begin{table}[H]\centering
\caption{Example reasoning prompts (native script)}
\small
\begin{tabular}{p{1.3cm}p{1cm}p{7cm}p{1.8cm}}\toprule
Lang. & Task & Prompt & Answer\\\midrule
Hindi & Direct & \deva{अगर राम की उम्र 55 वर्ष हो और सीता की उम्र 66 वर्ष हो, तो किसकी उम्र ज़्यादा होगी? उत्तर}: & \deva{सीता}\\
Hindi & Multihop & \deva{राम की उम्र, सीता की उम्र से कम है। सीता की उम्र, सुरेश की उम्र से कम है। राम और सुरेश में से किसकी उम्र ज़्यादा है? उत्तर}: & \deva{सुरेश}\\
Nepali & Direct & \deva{यदि विकासको उमेर 23 वर्ष छ र सीताको उमेर 23 वर्ष छ भने, कसको उमेर बढी होला? जवाफ}: & \deva{दुवै बराबर छन्}\\
Nepali & Multihop & \deva{चश्माको मूल्य, कलमको मूल्य भन्दा कम छ। कलमको मूल्य, पङ्खाको मूल्य भन्दा कम छ। चश्मा र पङ्खा मध्ये कसको मूल्य बढी छ? जवाफ}: & \deva{पङ्खा}\\
\bottomrule\end{tabular}
\end{table}
\subsection{Finetuning}
Finetuning started from each language's own pretrained checkpoint
(weights only -- optimizer/scheduler state was reinitialized fresh, and
the step counter restarts at 0). Each example is tokenized as
$\text{prompt}+\text{answer}+\texttt{EOS}$, and the loss is masked so only
the answer span (not the prompt) is supervised. Unlike pretraining, each
step draws one batch of 32 examples directly (no gradient accumulation),
sampled with replacement from the 6,000-example train set -- the same
resumable-by-step-counter philosophy as pretraining.

\begin{table}[H]\centering
\caption{Finetuning hyperparameters}
\small
\begin{tabular}{ll}\toprule
Setting & Value\\\midrule
Optimizer & AdamW, $\beta=(0.9,0.95)$\\
Peak learning rate & $2\times10^{-5}$\\
Schedule & linear warmup (50 steps) $\to$ cosine decay to 10\% of peak\\
Batch size & 32 (no gradient accumulation)\\
Max sequence length & 96 tokens\\
Gradient clipping & max norm 1.0\\
Loss & masked cross-entropy, answer span only\\
Steps run & 3,000 (Hindi) / 5,000 (Nepali)\\
\bottomrule\end{tabular}
\end{table}

\begin{figure}[H]\centering
\includegraphics[width=\textwidth]{finetune_loss_curves.png}
\caption{Reasoning finetuning train/validation loss (masked cross-entropy,
answer span only). Hindi converges smoothly; Nepali shows a plateau
followed by a step-change drop around step $\sim$1,900, plausibly from the
cosine-decay schedule being recomputed against a new \texttt{max\_steps}
each time the run was resumed with an extended step budget.}
\label{fig:finetune-loss}
\end{figure}


Nepali was extended from 1,000 (config default) to 3,000, then to 5,000
steps specifically to test whether its accuracy gap versus Hindi was a
step-budget artifact or a real ceiling

\subsection{Reasoning Results}

\begin{table}[H]\centering
\caption{Answer-restricted accuracy: pretrained vs.\ finetuned (same held-out test set)}
\small
\begin{tabular}{lrrrr}\toprule
& \multicolumn{2}{c}{Hindi} & \multicolumn{2}{c}{Nepali}\\
Task & Pretrained & Finetuned & Pretrained & Finetuned\\\midrule
Overall & \textbf{32.1\%} & 22.0\% & \textbf{36.6\%} & 11.9\%\\
Direct & 36.7\% & 11.7\% & 40.8\% & 11.7\%\\
Transitive & 25.1\% & 26.4\% & 30.8\% & \textbf{0.9\%}\\
Multihop & 31.7\% & 35.6\% & 35.6\% & 24.4\%\\
\bottomrule\end{tabular}
\end{table}


The pretrained models achieved 0\% exact-match accuracy. Finetuning
substantially improved performance, with a larger improvement for
Hindi.

Increasing Nepali finetuning from 3,000 to 5,000 steps did not close
the gap, suggesting that insufficient optimization steps alone do not
explain the difference.

\subsection{Failure Analysis}
Hindi errors generally remained within the entities appearing in the
prompt. Nepali sometimes generated entity names absent from the
current prompt but present in the training pool, indicating a possible
entity-binding or grounding problem.

Attention statistics also changed after finetuning, with overall
entropy decreasing in both models. These changes show that finetuning
altered attention organization, although they do not by themselves
establish the cause of the reasoning gap.

\section{Examples:}
\begin{table}[H]\centering
\caption{Example prompts: pretrained vs.\ finetuned predictions (real model outputs)}
\small
\begin{tabular}{p{1.3cm}p{1cm}p{5.2cm}p{1.6cm}p{2.6cm}p{1.6cm}}\toprule
Lang.\ & Task & Prompt & Answer & Pretrained pred. & Finetuned pred.\\\midrule
Model H (Hindi) & Multihop & \deva{राम की उम्र, सीता की उम्र से कम है। सीता की उम्र, सुरेश की उम्र से कम है। राम और सुरेश में से किसकी उम्र ज़्यादा है? उत्तर}: & \deva{सुरेश} & \deva{राम की उम्र, राम की उम्र}, & \textbf{\deva{सुरेश}}\\
Model H (Hindi) & Transitive & \deva{पता है कि राम की ऊँचाई, विकास की ऊँचाई से कम है, और विकास की ऊँचाई, सीता की ऊँचाई से कम है। इनमें सबसे कम ऊँचाई किसकी है, बताइए। उत्तर}: & \deva{राम} & - "\deva{अब तक, राम की} & \textbf{\deva{राम}}\\
Model H (Hindi) & Direct & \deva{अगर विकास की ऊँचाई 160 सेंटीमीटर हो और सुनीता की ऊँचाई 108 सेंटीमीटर हो, तो किसकी ऊँचाई कम होगी? उत्तर}: & \deva{सुनीता} & -\deva{पश्चिमी भाग में, यह एक} & \deva{दोनों बराबर हैं}\\
Model L (Nepali) & Multihop & \deva{चश्माको मूल्य, कलमको मूल्य भन्दा कम छ। कलमको मूल्य, पङ्खाको मूल्य भन्दा कम छ। चश्मा र पङ्खा मध्ये कसको मूल्य बढी छ? जवाफ}: & \deva{पङ्खा} & "\deva{अनलाइन" को मूल्य}, & \textbf{\deva{पङ्खा}}\\
Model L (Nepali) & Direct & \deva{यदि विकासको उमेर 23 वर्ष छ र सीताको उमेर 23 वर्ष छ भने, कसको उमेर बढी होला? जवाफ}: & \deva{दुवै बराबर छन्} & \textit{(empty)} & \textbf{\deva{दुवै बराबर छन्}}\\
Model L (Nepali) & Direct & \deva{यदि मनोजको उचाइ 160 सेन्टिमिटर छ र सुनिताको उचाइ 108 सेन्टिमिटर छ भने, कसको उचाइ कम होला? जवाफ}: & \deva{सुनिता} & \textit{(empty)} & \deva{सुनिल}\\
\bottomrule\end{tabular}
\end{table}


\section{Conclusion}
Two independent decoder-only Transformer language models were
implemented and trained from scratch for Hindi and Nepali. Both used
approximately 500M-token corpora, 8,000-token vocabularies, and the
same approximately 25M-parameter architecture.

The experiments show that intrinsic language modeling, generation,
attention organization, and reasoning transfer measure different
properties. Hindi obtained 36.75\% reasoning accuracy after
finetuning, compared with 8.75\% for Nepali, while intrinsic metrics
showed a more nuanced picture because of tokenizer differences.

\end{document}
